% abntex2-config.tex - Default abnTeX2 configuration
% Get Thesis Done (GTD)
%
% This file provides common abnTeX2 customizations for Brazilian academic theses.
% It can be \input{} from main.tex if desired, but main.tex works without it.
%
% To use: add % abntex2-config.tex - Default abnTeX2 configuration
% Get Thesis Done (GTD)
%
% This file provides common abnTeX2 customizations for Brazilian academic theses.
% It can be \input{} from main.tex if desired, but main.tex works without it.
%
% To use: add % abntex2-config.tex - Default abnTeX2 configuration
% Get Thesis Done (GTD)
%
% This file provides common abnTeX2 customizations for Brazilian academic theses.
% It can be \input{} from main.tex if desired, but main.tex works without it.
%
% To use: add % abntex2-config.tex - Default abnTeX2 configuration
% Get Thesis Done (GTD)
%
% This file provides common abnTeX2 customizations for Brazilian academic theses.
% It can be \input{} from main.tex if desired, but main.tex works without it.
%
% To use: add \input{abntex2-config} in main.tex preamble (after \documentclass).
% To customize: copy this file to your project root and modify as needed.
%
% NOTE: This is the default configuration. Users with specific university templates
% should create a custom configuration file or replace this one entirely.

% --- Spacing ---
% abnTeX2 defaults to 1.5 line spacing (ABNT NBR 14724 requirement).
% Uncomment to override:
% \OnehalfSpacing  % 1.5 spacing (default, ABNT compliant)
% \SingleSpacing   % Single spacing (for block quotes, footnotes)

% --- Paragraph indentation ---
% ABNT NBR 14724 requires 1.25cm paragraph indentation.
\setlength{\parindent}{1.25cm}

% --- Paragraph spacing ---
% ABNT NBR 14724: no extra space between paragraphs.
\setlength{\parskip}{0cm}

% --- Header and footer ---
% abnTeX2 default: page number in upper right corner (ABNT compliant).
% Uncomment to customize:
% \makepagestyle{abntheadings}
% \makeevenhead{abntheadings}{\ABNTEXfontereduzida\thepage}{}{}
% \makeoddhead{abntheadings}{}{}{\ABNTEXfontereduzida\thepage}
% \makeheadrule{abntheadings}{\textwidth}{\normalrulethickness}

% --- Caption formatting ---
% ABNT requires captions above figures/tables with specific formatting.
% abnTeX2 handles this by default. Uncomment to customize:
% \usepackage{caption}
% \captionsetup{
%   font=small,
%   labelfont=bf,
%   labelsep=endash,
%   justification=centering,
% }

% --- Float placement ---
% Prevent figures and tables from floating too far from their reference.
% \usepackage{placeins}  % Provides \FloatBarrier

% --- Code listings ---
% Uncomment if the thesis includes code listings.
% \usepackage{listings}
% \lstset{
%   basicstyle=\ttfamily\small,
%   breaklines=true,
%   frame=single,
%   numbers=left,
%   numberstyle=\tiny,
% }

% --- Mathematical packages ---
% Uncomment if the thesis uses mathematical notation.
% \usepackage{amsmath}
% \usepackage{amssymb}
% \usepackage{amsthm}

% --- Color ---
% Uncomment for colored text or hyperlinks.
% \usepackage{xcolor}
% \hypersetup{
%   colorlinks=true,
%   linkcolor=blue,
%   citecolor=blue,
%   urlcolor=blue,
% }

% --- Glossary / Acronyms ---
% Uncomment if the thesis uses a formal glossary or acronym list.
% \usepackage[acronym,nonumberlist]{glossaries}
% \makeglossaries
 in main.tex preamble (after \documentclass).
% To customize: copy this file to your project root and modify as needed.
%
% NOTE: This is the default configuration. Users with specific university templates
% should create a custom configuration file or replace this one entirely.

% --- Spacing ---
% abnTeX2 defaults to 1.5 line spacing (ABNT NBR 14724 requirement).
% Uncomment to override:
% \OnehalfSpacing  % 1.5 spacing (default, ABNT compliant)
% \SingleSpacing   % Single spacing (for block quotes, footnotes)

% --- Paragraph indentation ---
% ABNT NBR 14724 requires 1.25cm paragraph indentation.
\setlength{\parindent}{1.25cm}

% --- Paragraph spacing ---
% ABNT NBR 14724: no extra space between paragraphs.
\setlength{\parskip}{0cm}

% --- Header and footer ---
% abnTeX2 default: page number in upper right corner (ABNT compliant).
% Uncomment to customize:
% \makepagestyle{abntheadings}
% \makeevenhead{abntheadings}{\ABNTEXfontereduzida\thepage}{}{}
% \makeoddhead{abntheadings}{}{}{\ABNTEXfontereduzida\thepage}
% \makeheadrule{abntheadings}{\textwidth}{\normalrulethickness}

% --- Caption formatting ---
% ABNT requires captions above figures/tables with specific formatting.
% abnTeX2 handles this by default. Uncomment to customize:
% \usepackage{caption}
% \captionsetup{
%   font=small,
%   labelfont=bf,
%   labelsep=endash,
%   justification=centering,
% }

% --- Float placement ---
% Prevent figures and tables from floating too far from their reference.
% \usepackage{placeins}  % Provides \FloatBarrier

% --- Code listings ---
% Uncomment if the thesis includes code listings.
% \usepackage{listings}
% \lstset{
%   basicstyle=\ttfamily\small,
%   breaklines=true,
%   frame=single,
%   numbers=left,
%   numberstyle=\tiny,
% }

% --- Mathematical packages ---
% Uncomment if the thesis uses mathematical notation.
% \usepackage{amsmath}
% \usepackage{amssymb}
% \usepackage{amsthm}

% --- Color ---
% Uncomment for colored text or hyperlinks.
% \usepackage{xcolor}
% \hypersetup{
%   colorlinks=true,
%   linkcolor=blue,
%   citecolor=blue,
%   urlcolor=blue,
% }

% --- Glossary / Acronyms ---
% Uncomment if the thesis uses a formal glossary or acronym list.
% \usepackage[acronym,nonumberlist]{glossaries}
% \makeglossaries
 in main.tex preamble (after \documentclass).
% To customize: copy this file to your project root and modify as needed.
%
% NOTE: This is the default configuration. Users with specific university templates
% should create a custom configuration file or replace this one entirely.

% --- Spacing ---
% abnTeX2 defaults to 1.5 line spacing (ABNT NBR 14724 requirement).
% Uncomment to override:
% \OnehalfSpacing  % 1.5 spacing (default, ABNT compliant)
% \SingleSpacing   % Single spacing (for block quotes, footnotes)

% --- Paragraph indentation ---
% ABNT NBR 14724 requires 1.25cm paragraph indentation.
\setlength{\parindent}{1.25cm}

% --- Paragraph spacing ---
% ABNT NBR 14724: no extra space between paragraphs.
\setlength{\parskip}{0cm}

% --- Header and footer ---
% abnTeX2 default: page number in upper right corner (ABNT compliant).
% Uncomment to customize:
% \makepagestyle{abntheadings}
% \makeevenhead{abntheadings}{\ABNTEXfontereduzida\thepage}{}{}
% \makeoddhead{abntheadings}{}{}{\ABNTEXfontereduzida\thepage}
% \makeheadrule{abntheadings}{\textwidth}{\normalrulethickness}

% --- Caption formatting ---
% ABNT requires captions above figures/tables with specific formatting.
% abnTeX2 handles this by default. Uncomment to customize:
% \usepackage{caption}
% \captionsetup{
%   font=small,
%   labelfont=bf,
%   labelsep=endash,
%   justification=centering,
% }

% --- Float placement ---
% Prevent figures and tables from floating too far from their reference.
% \usepackage{placeins}  % Provides \FloatBarrier

% --- Code listings ---
% Uncomment if the thesis includes code listings.
% \usepackage{listings}
% \lstset{
%   basicstyle=\ttfamily\small,
%   breaklines=true,
%   frame=single,
%   numbers=left,
%   numberstyle=\tiny,
% }

% --- Mathematical packages ---
% Uncomment if the thesis uses mathematical notation.
% \usepackage{amsmath}
% \usepackage{amssymb}
% \usepackage{amsthm}

% --- Color ---
% Uncomment for colored text or hyperlinks.
% \usepackage{xcolor}
% \hypersetup{
%   colorlinks=true,
%   linkcolor=blue,
%   citecolor=blue,
%   urlcolor=blue,
% }

% --- Glossary / Acronyms ---
% Uncomment if the thesis uses a formal glossary or acronym list.
% \usepackage[acronym,nonumberlist]{glossaries}
% \makeglossaries
 in main.tex preamble (after \documentclass).
% To customize: copy this file to your project root and modify as needed.
%
% NOTE: This is the default configuration. Users with specific university templates
% should create a custom configuration file or replace this one entirely.

% --- Spacing ---
% abnTeX2 defaults to 1.5 line spacing (ABNT NBR 14724 requirement).
% Uncomment to override:
% \OnehalfSpacing  % 1.5 spacing (default, ABNT compliant)
% \SingleSpacing   % Single spacing (for block quotes, footnotes)

% --- Paragraph indentation ---
% ABNT NBR 14724 requires 1.25cm paragraph indentation.
\setlength{\parindent}{1.25cm}

% --- Paragraph spacing ---
% ABNT NBR 14724: no extra space between paragraphs.
\setlength{\parskip}{0cm}

% --- Header and footer ---
% abnTeX2 default: page number in upper right corner (ABNT compliant).
% Uncomment to customize:
% \makepagestyle{abntheadings}
% \makeevenhead{abntheadings}{\ABNTEXfontereduzida\thepage}{}{}
% \makeoddhead{abntheadings}{}{}{\ABNTEXfontereduzida\thepage}
% \makeheadrule{abntheadings}{\textwidth}{\normalrulethickness}

% --- Caption formatting ---
% ABNT requires captions above figures/tables with specific formatting.
% abnTeX2 handles this by default. Uncomment to customize:
% \usepackage{caption}
% \captionsetup{
%   font=small,
%   labelfont=bf,
%   labelsep=endash,
%   justification=centering,
% }

% --- Float placement ---
% Prevent figures and tables from floating too far from their reference.
% \usepackage{placeins}  % Provides \FloatBarrier

% --- Code listings ---
% Uncomment if the thesis includes code listings.
% \usepackage{listings}
% \lstset{
%   basicstyle=\ttfamily\small,
%   breaklines=true,
%   frame=single,
%   numbers=left,
%   numberstyle=\tiny,
% }

% --- Mathematical packages ---
% Uncomment if the thesis uses mathematical notation.
% \usepackage{amsmath}
% \usepackage{amssymb}
% \usepackage{amsthm}

% --- Color ---
% Uncomment for colored text or hyperlinks.
% \usepackage{xcolor}
% \hypersetup{
%   colorlinks=true,
%   linkcolor=blue,
%   citecolor=blue,
%   urlcolor=blue,
% }

% --- Glossary / Acronyms ---
% Uncomment if the thesis uses a formal glossary or acronym list.
% \usepackage[acronym,nonumberlist]{glossaries}
% \makeglossaries
